% CN 7.19 Exercises (p.274--280)

% ======================================================================
\begin{exercise}{1}
    \begin{enumerate}[label=(\alph*)]
        \item Write the statement that integer $d$ is a divisor of integer $n$ as
              a predicate logic expression, using integer multiplication and
              without using $\mid$.
        \item Write the statement that integer $n$ is prime as a predicate logic
              expression. You may use $\mid$ and $\le$.
        \item Write the statement that $d$ is the greatest common divisor of $m$
              and $n$ as a predicate logic expression, using $\mid$ and $\le$.
    \end{enumerate}
\end{exercise}

\begin{exercise}{2}
    A positive integer $n$ is \emph{perfect} if it equals the sum of its own
    proper divisors.
    \begin{enumerate}[label=(\alph*)]
        \item Show by direct calculation that the first two perfect numbers are
              6 and 28.
        \item Verify that 496 is perfect.
    \end{enumerate}
\end{exercise}

\begin{exercise}{3}
    Represent days of the week by $\Z_7$ (Sunday $=0$, Monday $=1$, \ldots). Let
    $d\in\Z_7$ be the day of the week of a given date this year.
    \begin{enumerate}[label=(\alph*)]
        \item Give an expression in $d$, using mod, for the day of the week of
              the same date next year.
        \item How would you modify it if there is a leap day (29 February)
              between that date and the same date next year?
        \item Let $d$ be the day of the week of 29 February in some leap year.
              Assuming the next leap year is four years later, give an expression
              for the weekday of the next 29 February.
    \end{enumerate}
\end{exercise}

\begin{exercise}{4}
    In (2.1) on p.54 we defined
    $\mathrm{LCG}(x)$ = the last 31 bits of $1103515245x+12345$. Restate this
    definition using the mod operation.
\end{exercise}

\begin{exercise}{5}
    The online game Primel (\url{https://converged.yt/primel/}) is about
    guessing an unknown five-digit prime (like Wordle). Non-prime guesses are not
    penalised but give no information; you have at most six prime guesses.
    \begin{enumerate}[label=(\alph*)]
        \item Are there any prohibitions on certain digits in certain positions
              in prime numbers in general? What does this mean for the digits in
              your first guess?
        \item The density of primes decreases as numbers get higher (e.g.\ more
              five-digit primes start with 1 than with 5). What five digits
              would be best in your first guess? Which five-digit prime do you
              recommend?
        \item Try asking ChatGPT for help choosing your next guess, and check
              whether its suggestions are actually prime. How many attempts are
              necessary before it gives an actual prime satisfying the conditions?
    \end{enumerate}
\end{exercise}

\begin{exercise}{6}
    \emph{A positive integer is a leap year iff it is a multiple of 4, unless it
    is also a multiple of 100, in which case it is not a leap year unless it is a
    multiple of 400.} Define the unary predicate $\mathrm{Leap}$ on $\N$. Using
    $\mathrm{Leap}$ and $\mid$, write this rule in predicate logic.
\end{exercise}

\begin{exercise}{7}
    Prove that the square of an odd number is odd and the square of an even
    number is even.
\end{exercise}

\begin{exercise}{8}
    Restate the definitions of the Caesar slide encryption and decryption
    functions (Exercise 2.15) using the mod operation.
\end{exercise}

% ======================================================================
\begin{exercise}{9}
    Prove each of the following.
    \begin{enumerate}[label=(\alph*)]
        \item $4\mid n$ iff the last two decimal digits of $n$ give a multiple
              of 4.
        \item $8\mid n$ iff the last three decimal digits of $n$ give a multiple
              of 8.
        \item $25\mid n$ iff the last two decimal digits of $n$ give a multiple
              of 25.
        \item Extend (a) and (b) to higher powers of 2.
        \item (Challenge) For which $d$ can divisibility of $n$ by $d$ be tested
              by looking only at some constant number of digits at the end of
              $n$?
    \end{enumerate}
\end{exercise}

\begin{exercise}{10}
    The \emph{digital sum} $\mathrm{ds}(n)$ of $n$ is the sum of its decimal
    digits, e.g.\ $\mathrm{ds}(1984)=1+9+8+4=22$.
    \begin{enumerate}[label=(\alph*)]
        \item Prove that for all $n\in\N$,
              $n\bmod3=\mathrm{ds}(n)\bmod3$.
        \item A similar method works for remainders mod 9. Explain why, briefly.
        \item The \emph{alternating digital sum} $\mathrm{ads}(n)$ alternately
              adds and subtracts digits, starting with addition at the right-hand
              end, e.g.\ $\mathrm{ads}(1984)=-1+9-8+4=4$. Prove that for all
              $n\in\N$, $n\bmod11=\mathrm{ads}(n)\bmod11$.
        \item Devise a similar technique for working out $n\bmod3$ from the bits
              of the binary representation of $n$.
    \end{enumerate}
\end{exercise}

\begin{exercise}{11}
    For each equation, use the method of Exercise 10(b) to work out by hand the
    remainder mod 9 of each side. Comment on what comparing them tells you.
    \begin{enumerate}[label=(\alph*)]
        \item $92\times31=2847$
        \item $92\times31=2852$
        \item $92\times31=2861$
        \item $2718281828\times3141592653=8539734219628209634$
    \end{enumerate}
\end{exercise}

% ======================================================================
\begin{exercise}{12}
    Use the Euclidean algorithm to find the gcd of the consecutive Fibonacci
    numbers 34 and 55. What do you notice about the other numbers found at each
    step?

    Prove, by induction on $n$, that for all $n\ge3$ the Euclidean algorithm uses
    $n-2$ subtractions to compute $\gcd(f_n,f_{n+1})$, and that every
    intermediate number found is also a Fibonacci number.
\end{exercise}

\begin{exercise}{13}
    Use the Extended Euclidean Algorithm to show that 86 and 99 are coprime, to
    express 1 as an integer linear combination of them, and to find the inverse
    of 86 in $\Z_{99}$.
\end{exercise}

\begin{exercise}{14}
    \begin{enumerate}[label=(\alph*)]
        \item Find $\varphi(n)$ for all $n$ up to 20.
        \item Find primitive roots for a selection of values of $n$.
    \end{enumerate}
\end{exercise}

\begin{exercise}{15}
    \textbf{Modular multiplication cryptosystem.} Messages and cyphertexts are
    strings over the 26-letter alphabet, with $\mathtt a,\dots,\mathtt z$ as
    $0,\dots,25\in\Z_{26}$. The key is $k\in\Z_{26}$, and
    \[
        c_i = k\,m_i \bmod 26.
    \]
    \begin{enumerate}[label=(\alph*)]
        \item Not all keys work properly. What restriction must be placed on
              members of $\Z_{26}$ so they work as keys? Specify the keyspace in
              mathematical notation.
        \item Define the decryption function.
        \item How would modular multiplication work for an arbitrary alphabet
              size $n$? Define its keyspace, encryption and decryption functions.
    \end{enumerate}
\end{exercise}

\begin{exercise}{16}
    Compute by hand, showing your working,
    \[
        1752^{31415926535897932384626433832795028841971693993751058209749445923078164}
        \bmod 125.
    \]
    (The exponent is $\lfloor10^{70}\pi\rfloor$.)
\end{exercise}

\begin{exercise}{17}
    Investigate the primitive roots of 11.
    \begin{enumerate}[label=(\alph*)]
        \item If $a$ is a primitive root of 11, what is the least $k$ such that
              $a^k\equiv1\pmod{11}$?
        \item Show that 2 is a primitive root of 11 by working out its powers
              $2^i$ (multiplying by 2 and reducing mod 11 as you go).
        \item Among the exponents $i$ in your list from (b), which are coprime to
              $\varphi(11)$?
        \item Using (c), list all the primitive roots of 11.
    \end{enumerate}
\end{exercise}

% ======================================================================
\begin{exercise}{18}
    Alice and Bob use Diffie-Hellman with public prime $p=11$ and primitive root
    $a=7$. Their public numbers are $y_A=2$, $y_B=8$. As the cryptanalyst, find
    their private numbers and the shared key.
\end{exercise}

\begin{exercise}{19}
    You have been using Diffie-Hellman with $p=17$ and $a=7$. Your General
    Manager decides that ``bigger is better'' and you will now use $p=18$ (with
    appropriate $a$).
    \begin{enumerate}[label=(\alph*)]
        \item What value of $a$ would you choose?
        \item What are the security implications of the change? Estimate e.g.\
              the percentage change in the time for an exhaustive attack.
    \end{enumerate}
\end{exercise}

\begin{exercise}{20}
    Explain how the security of Diffie-Hellman is affected if user A selects
    private key $x_A=p-1$ (where $p$ is the prime modulus).
\end{exercise}

\begin{exercise}{21}
    Devise a variant of Diffie-Hellman enabling three people A, B, C to arrive at
    a common secret key, such that:
    \begin{itemize}
        \item they initially have no secret information in common;
        \item the common key itself is not sent anywhere by anyone;
        \item it is hard for an eavesdropper (knowing only public information)
              to determine the key;
        \item no one or two of A, B, C can find the key without the cooperation
              of the other(s).
    \end{itemize}
\end{exercise}
